% from ega3-1-fr.tex
\oldpage[0\textsubscript{III}]{31}
    dès que $p-r+1+j_0<n$, donc
    $B_r^{pq}=B_\infty(E_2^{pq})$ pour $r<j_0-q+1$~; d'autre part,
    $\mathrm H^n(F_{\mathrm I}^p(K^{\bullet\bullet}))
    =\mathrm H^n(K^{\bullet\bullet})$ pour $p+j_0<n-1$.
\end{enumerate}
\end{proof}

\begin{env}[11.3.4]
\label{0.11.3.4-fr}
Supposons que le bicomplexe $K^{\bullet\bullet}=(K^{ij})$ soit tel que
$K^{ij}=0$ pour $i<0$ ou $j<0$. On sait qu'on peut alors définir pour tout
$p\in\mathbb Z$ un «~edge-homomorphisme~» canonique
\[
\label{0.11.3.4.1-fr}
  'E_2^{p0}(K^{\bullet\bullet})\longrightarrow
  \mathrm H^p(K^{\bullet\bullet})
  \tag{11.3.4.1}
\]
(M, XV, 6). Rappelons rapidement que cela est dû, d'une part à ce que l'on a
dans la suite spectrale $'E(K^{\bullet\bullet})$,
$Z_r^{p0}=Z_{\mathrm I}^p(Z_{\mathrm{II}}^0(K^{\bullet\bullet}))$ pour
$2\leq r\leq+\infty$, et de l'autre au fait que
$\mathrm H^p(F_{\mathrm I}^{p+1}(K^{\bullet\bullet}))=0$, si bien que
l'isomorphisme
$\beta^{p0}:'E_\infty^{p0}\xrightarrow{\sim}
\mathrm H^p(F_{\mathrm I}^p)/\mathrm H^p(F_{\mathrm I}^{p+1})$ donne un
homomorphisme
$'E_\infty^{p0}\to\mathrm H^p(F_{\mathrm I}^p(K^{\bullet\bullet}))
\to\mathrm H^p(K^{\bullet\bullet})$~; l'égalité de tous les $Z_r^{p0}$ permet
alors de définir des homomorphismes canoniques
$'E_r^{p0}\to'E_s^{p0}$ pour $r\leq s$, et en particulier un homomorphisme
$'E_2^{p0}\to'E_\infty^{p0}$, d'où par composition l'edge-homomorphisme
$'E_2^{p0}\to\mathrm H^p(K^{\bullet\bullet})$~; en outre, on vérifie aussitôt
que, à la classe mod. $B_2^{p0}$ d'un élément
$z\in Z_{\mathrm{II}}^0(K^{\bullet\bullet})\subset K^{p0}$ tel que $d'z=0$,
l'edge-homomorphisme ainsi défini fait correspondre dans $'E_\infty^{p0}$ la
classe de $z$ mod. $B_\infty^{p0}$, puis à cette dernière la classe de
cohomologie de $z$ dans $\mathrm H^p(K^{\bullet\bullet})$. On voit donc
finalement que \emph{l'edge-homomorphisme
(\hyperref[0.11.3.4.1-fr]{11.3.4.1}) provient, par passage à la cohomologie,
de l'injection canonique}
$Z_{\mathrm{II}}^0(K^{\bullet\bullet})\to K^{\bullet\bullet}$ (où
$K^{\bullet\bullet}$ est considéré comme complexe simple). On interprète
naturellement de même l'edge-homomorphisme
\[
\label{0.11.3.4.2-fr}
  ''E_2^{p0}(K^{\bullet\bullet})\longrightarrow
  \mathrm H^p(K^{\bullet\bullet})
  \tag{11.3.4.2}
\]
comme provenant de l'injection canonique
$Z_{\mathrm I}^0(K^{\bullet\bullet})\to K^{\bullet\bullet}$.
\end{env}

\begin{env}[11.3.5]
\label{0.11.3.5-fr}
Soit maintenant $K_{\bullet\bullet}=(K_{ij})$ un bicomplexe de
$\mathcal C$ dont les deux opérateurs de dérivation sont de degré $-1$. Nous
écrirons alors $K_{i,\bullet}$ (resp. $K_{\bullet,j}$) le complexe simple
$(K_{ij})_{j\in\mathbb Z}$ (resp. $(K_{ij})_{i\in\mathbb Z}$),
$\mathrm H_p^{\mathrm{II}}(K_{i,\bullet})$ (resp.
$\mathrm H_p^{\mathrm I}(K_{\bullet,j})$) son $p$-ème objet d'homologie,
$\mathrm H_p^{\mathrm{II}}(K_{\bullet\bullet})$ (resp.
$\mathrm H_p^{\mathrm I}(K_{\bullet\bullet})$) le complexe
$(\mathrm H_p^{\mathrm{II}}(K_{i,\bullet}))_{i\in\mathbb Z}$ (resp.
$(\mathrm H_p^{\mathrm I}(K_{\bullet,j}))_{j\in\mathbb Z}$),
$\mathrm H_q^{\mathrm I}(\mathrm H_p^{\mathrm{II}}(K_{\bullet\bullet}))$
(resp.
$\mathrm H_q^{\mathrm{II}}(\mathrm H_p^{\mathrm I}(K_{\bullet\bullet}))$)
son $q$-ème objet d'homologie~; notations analogues pour les objets des cycles
et objets des bords~; enfin, $\mathrm H_n(K_{\bullet\bullet})$ désignera
(lorsqu'il existe) le $n$-ème objet d'homologie du complexe simple (à
opérateur de dérivation de degré $-1$) défini par $K_{\bullet\bullet}$.

Soit $K^{\prime\bullet\bullet}=(K^{\prime ij})$ avec
$K^{\prime ij}=K_{-i,-j}$ le bicomplexe à opérateurs de dérivation de degrés
$+1$ associé à $K_{\bullet\bullet}$~; par définition, les
\emph{suites spectrales} de $K_{\bullet\bullet}$ sont celles de
$K^{\prime\bullet\bullet}$, que l'on écrit $'E(K_{\bullet\bullet})$ et
$''E(K_{\bullet\bullet})$, où l'on change toutefois les notations, en posant
\[
  'E_{pq}^r(K_{\bullet\bullet})
  ='E_r^{-p,-q}(K^{\prime\bullet\bullet}),\qquad
  ''E_{pq}^r(K_{\bullet\bullet})
  =''E_r^{-p,-q}(K^{\prime\bullet\bullet})
\]
pour $2\leq r\leq\infty$. Avec ces notations, on a
\[
  'E_{pq}^2(K_{\bullet\bullet})
  =\mathrm H_p^{\mathrm I}(\mathrm H_q^{\mathrm{II}}(K_{\bullet\bullet})),
  \qquad
  ''E_{pq}^2(K_{\bullet\bullet})
  =\mathrm H_q^{\mathrm{II}}(\mathrm H_p^{\mathrm I}(K_{\bullet\bullet})).
\]
Pour éviter des erreurs de signe, il sera en général préférable, pour les
relations entre ces suites spectrales et leur aboutissement, de revenir au
complexe $K^{\prime\bullet\bullet}$. Notons toutefois les critères
correspondant à (\hyperref[0.11.3.3-fr]{11.3.3})~:
\end{env}

\begin{env}[11.3.6]
\label{0.11.3.6-fr}
Les suites spectrales $'E(K_{\bullet\bullet})$ et
$''E(K_{\bullet\bullet})$ sont \emph{birégulières} dans les cas suivants~:
\begin{enumerate}
  \item[(a)] Il existe $i_0$ et $j_0$ tels que $K_{ij}=0$ pour $i>i_0$ ou
    pour $j>j_0$ (resp. pour $i<i_0$
    